\chapter{Hash-Based Signatures: One-Time and Few-Time}
\label{ch:hashots}

\section{Why this chapter matters}
The previous part built signatures on lattice problems. This part builds them on a completely different foundation -- nothing but a cryptographic hash function. There is no number theory, no lattice, no algebraic structure for an attacker to exploit: if the hash is secure, so is the signature. That makes hash-based signatures the most conservative post-quantum option, and it is why NIST standardized one (SLH-DSA, FIPS 205) as a hedge against a future structural break in lattices.

Because the foundation is so different, we introduce it the same way we introduced lattices: from the ground up. This chapter builds the atomic pieces -- one-time and few-time signatures -- and the next assembles them into a many-time, stateless scheme with Merkle trees. Only then, in Chapter~\ref{ch:slhdsa}, do we turn to the FIPS 205 standard itself.

\section{Learning goals}
After reading this chapter, you should be able to:

\begin{itemize}
  \item state the three security properties a hash function must provide, and why Grover only mildly weakens them;
  \item construct and verify a Lamport one-time signature, and explain why it is strictly one-time;
  \item explain how Winternitz (WOTS\textsuperscript{+}) trades hashing time for much shorter signatures;
  \item describe the few-time signature FORS and how its security bound changes with the signature count.
\end{itemize}

\section{What a secure hash function guarantees}
Earlier chapters used hash functions such as SHAKE only as \emph{expanders} -- deterministic ways to stretch a short seed into a matrix or noise (Chapters~\ref{ch:cbd}--\ref{ch:fips203}). Here a hash is the \emph{sole source of security}, so we state precisely what is assumed of $H:\{0,1\}^*\to\{0,1\}^n$.

\begin{definition}[Hash-function security]
A cryptographic hash $H$ is assumed to be
\begin{itemize}
  \item \emph{preimage resistant} (one-way): given $y$, it is hard to find any $x$ with $H(x)=y$;
  \item \emph{second-preimage resistant}: given $x$, it is hard to find $x'\neq x$ with $H(x')=H(x)$;
  \item \emph{collision resistant}: it is hard to find any pair $x\neq x'$ with $H(x)=H(x')$.
\end{itemize}
\end{definition}

That is the whole assumption. Because there is no hidden algebraic structure, there is no Shor-style break; the only quantum help is Grover's search (Chapter~\ref{ch:quantum}), which cuts preimage resistance from $n$ bits to about $n/2$. Hash-based schemes simply use a larger output $n$ (16, 24, or 32 bytes) to keep a comfortable margin. The threat that shatters RSA leaves hashing merely needing bigger parameters.

\section{The Lamport one-time signature}
The simplest possible hash-based signature, due to Lamport~\cite{lamport1979}, already contains the key idea: \emph{reveal a hash preimage as proof}.

To sign an $m$-bit message, the signer prepares, for each bit position, \emph{two} secret random values -- one to reveal if that bit is $0$, one if it is $1$ -- and publishes their hashes. Signing a message reveals exactly one secret per position: the one matching the bit. A verifier hashes each revealed secret and checks it against the correct public value.

\begin{figure}[H]
\centering
\begin{tikzpicture}[every node/.style={font=\scriptsize}]
  \node (lbl0) at (0,1.4) {position $i$:};
  \node[flownode] (s0) at (2.7,1.4) {$sk_{i,0}$};
  \node[flownode] (s1) at (6.8,1.4) {$sk_{i,1}$};
  \node[keynode] (p0) at (2.7,0) {$pk_{i,0}=H(sk_{i,0})$};
  \node[keynode] (p1) at (6.8,0) {$pk_{i,1}=H(sk_{i,1})$};
  \draw[flowarrow] (s0) -- (p0);
  \draw[flowarrow] (s1) -- (p1);
  \node[pqcorange, align=center] at (4.75,-1.0)
    {if the message bit is $0$, reveal $sk_{i,0}$;\quad if it is $1$, reveal $sk_{i,1}$};
\end{tikzpicture}
\caption{One position of a Lamport key. Each of the $m$ message-bit positions has two secret values whose hashes are public. A signature reveals exactly one secret per position -- the one matching the message bit -- and the verifier confirms it hashes to the published value. The unrevealed secrets stay hidden by preimage resistance.}
\label{fig:lamport}
\end{figure}

\begin{algorithm}[H]
\caption{Lamport one-time signature}
\label{alg:lamport}
\begin{algorithmic}[1]
\Statex \textbf{KeyGen:}
\State for $i=0,\dots,m-1$ and $b\in\{0,1\}$, sample $sk_{i,b}\gets\{0,1\}^n$ and set $pk_{i,b}\gets H(sk_{i,b})$
\State \Return secret key $(sk_{i,b})$, public key $(pk_{i,b})$
\Statex \textbf{Sign} message $M=(M_0,\dots,M_{m-1})\in\{0,1\}^m$:
\State \Return $\sigma=(sk_{0,M_0},\ sk_{1,M_1},\ \dots,\ sk_{m-1,M_{m-1}})$
\Statex \textbf{Verify} $(M,\sigma=(\sigma_0,\dots,\sigma_{m-1}))$:
\State \Return \emph{accept} iff $H(\sigma_i)=pk_{i,M_i}$ for every $i$
\end{algorithmic}
\end{algorithm}

\begin{remark}[Why it is strictly one-time]
A single signature reveals one secret per position -- the ones matching $M$. To forge a \emph{different} message $M'$, an attacker would need $sk_{i,M'_i}$ at each position where $M'_i\neq M_i$, and those secrets are still protected by preimage resistance. But two signatures on different messages jointly expose \emph{both} secrets at every differing position, after which any message can be forged. Hence a Lamport key is safe for exactly one signature.
\end{remark}

Lamport keys and signatures are large (each is $2m$ or $m$ hash values). The next construction keeps the one-time security but shrinks the signature dramatically.

\section{Winternitz: trading time for size (WOTS\textsuperscript{+})}
Instead of two independent secrets per bit, Winternitz uses one \emph{hash chain} per base-$w$ digit of the message. Starting from a secret $sk$, repeated hashing produces a chain, and the public value is its end:
\[
sk \;\xrightarrow{H}\; H(sk) \;\xrightarrow{H}\; \cdots \;\xrightarrow{H}\; H^{w-1}(sk)=pk.
\]
To sign the digit $m$, the signer reveals the intermediate point $H^{m}(sk)$; the verifier hashes it the remaining $w-1-m$ times and checks it reaches $pk$.

\begin{figure}[H]
\centering
\begin{tikzpicture}[node distance=10mm, every node/.style={font=\scriptsize}]
  \node[flownode] (sk) {$sk$};
  \node[flownode, right=of sk] (h1) {$H(sk)$};
  \node[flownode, right=of h1] (h2) {$H^2(sk)$};
  \node[right=6mm of h2] (dots) {$\cdots$};
  \node[keynode, right=6mm of dots] (pk) {$pk=H^{w-1}(sk)$};
  \draw[flowarrow] (sk) -- (h1);
  \draw[flowarrow] (h1) -- (h2);
  \draw[flowarrow] (h2) -- (dots);
  \draw[flowarrow] (dots) -- (pk);
  \draw[decorate,decoration={brace,amplitude=4pt},pqcorange]
        (h2.north west) -- node[above=3pt]{signature $=H^{m}(sk)$} (h2.north east);
\end{tikzpicture}
\caption{One WOTS\textsuperscript{+} hash chain of length $w$. The secret is the start, the public value is the end, and a signature reveals an intermediate point. A verifier hashes forward to check it reaches the public end. A full message digest is split into base-$w$ digits, each signed by its own chain.}
\label{fig:wots-chain}
\end{figure}

There is a catch a checksum must fix. Revealing $H^m(sk)$ lets anyone hash \emph{forward} to $H^{m'}(sk)$ for any $m'>m$, so an attacker could ``advance'' a digit to a larger value for free. WOTS\textsuperscript{+} appends a short \emph{checksum} of the digits, itself signed by extra chains, arranged so that increasing any message digit forces \emph{decreasing} some checksum digit -- which would require moving a chain \emph{backward}, i.e.\ inverting $H$. The parameter $w$ tunes a trade-off: larger $w$ means fewer, longer chains (shorter signatures, more hashing), smaller $w$ the reverse. Like Lamport, a WOTS\textsuperscript{+} key is safe for exactly one signature.

\section{FORS: a few-time signature}
One-time keys are unforgiving: a single reuse is catastrophic. The stateless scheme of the next chapters uses \textbf{FORS} (Forest Of Random Subsets), a \emph{few-time} signature, inside a very large address space. This does not authorize key reuse; it makes the security analysis degrade more gradually under the tiny probability of an accidental address collision.

\begin{definition}[FORS]
FORS consists of $k$ independent Merkle trees, each of height $a$ (so each has $2^a$ leaves, and each leaf hides a secret value). A message digest is split into $k$ indices, one per tree; the signature reveals, for each tree, the secret at the selected leaf together with its authentication path to that tree's root. The $k$ roots are hashed together into a single FORS public key.
\end{definition}

Because each tree reveals only the \emph{one} leaf its index selects, the security bound degrades with the number of signatures and with any address collisions. SLH-DSA uses message randomization and a huge address space to make such collisions extremely unlikely. An implementation must never deliberately reuse a FORS or WOTS address; WOTS leaf reuse remains a serious event. (The Merkle trees FORS uses are the subject of the next chapter.)

\section{Summary}
This chapter built the atoms of hash-based signing:

\begin{itemize}
  \item security rests only on hash preimage/collision resistance, which Grover weakens merely to $n/2$ bits;
  \item \textbf{Lamport} signs once by revealing one preimage per message bit -- simple but large;
  \item \textbf{WOTS\textsuperscript{+}} replaces bit-pairs with hash chains plus a checksum, keeping one-time security at a fraction of the size;
  \item \textbf{FORS} is a \emph{few-time} signature whose security bound degrades with the number of signatures; address reuse is never an implementation strategy.
\end{itemize}

Every one of these keys is still usable only once (or a few times). The next chapter removes that limitation: a Merkle tree binds many one-time keys under a single public root, and a hypertree of such trees, addressed by hashing the message, yields a stateless many-time signature.

\printbibliography[heading=chapterrefs]
